\begin{abstract}
Influence maximization is dominated by reverse-reachability (RR) sampling, whose guarantees rest
on a coverage representation: for a random set $R$ independent of the seed set $S$, the spread is
$c \cdot \Pr[S \cap R \neq \emptyset]$, a non-decreasing submodular function of $S$. We study
targeted influence maximization under a threshold-dependent overexposure model in which a node
becomes positively active only when the weighted influence it receives from its in-neighbours lands
inside a node-specific window $[\theta^\kappa_v, \theta^\tau_v]$. We first delimit what is and is
not ruled out here. On a one-hop instance with two candidates of weight $0.4$ into a single target,
the model's own window distribution gives $F_D(\{a\}) = 0.48$ and $F_D(\{a,b\}) = 0.32$: the
objective is non-monotone, so no non-negative seed-independent coverage function can match it, and
the standard RR representation is unavailable. The statement is deliberately narrow --- it does not
exclude signed decompositions, restricted-instance algorithms, or Monte-Carlo estimation, and it
does not imply that learning must help. It also requires naming the threshold distribution: under
uniform-threshold linear thresholds the same instance is monotone, so the two readings cannot be
conflated. The practical question therefore becomes how to spend a simulation budget well. We
characterise the regime where that question bites: across nine graphs, saturation is controlled by
the seed \emph{fraction} $|S|/n$ rather than the absolute budget, and once $|S|/n \ge 20\%$ a
substantial share of candidates have a \emph{negative} marginal gain, so adding a seed can reduce
the target's activation probability. We then evaluate whether a learned marginal-gain predictor can
reduce the number of cascade simulations needed to reach a given quality, against a training-free
baseline and an adaptive Monte-Carlo allocation that uses the same verification machinery. We report
both what the evidence supports and what it does not.

\keywords{Influence maximization \and Overexposure \and Marginal gain estimation \and
Simulation budget \and Negative results}
\end{abstract}
